\renewcommand{\theposttitle}{How Progress Is Like Crossing a River}
\renewcommand{\thepostdate}{2025-04-29}
\vspace*{\fill}
\begin{center}
\begin{minipage}{0.92\textwidth}
\small
A quick and easy visual metaphor anyone can make ↴
\end{minipage}
\end{center}
\vspace{4mm}
\begin{center}
\includegraphics[width=\textwidth, height=0.55\textheight, keepaspectratio]{C:/Users/James/Documents/GitHub/VisualIdeasBook/data/images/b8098475b98eacf77cfe5cd2e20796e6.jpeg}
\end{center}
\vspace{2mm}
\begin{center}
\begin{minipage}{0.92\textwidth}
\small
A useful way to show the Pareto Principle (80/20 Rule). You can use this for anything where you want to highlight the importance of the smaller element—like the handful of customers who drive most of your revenue.

The big square is the majority (80\%), while that small detached piece represents the minority (20\%).

Know what else can show a majority/minority relationship? That’s right! A pie chart! So you can think of this framework as a variation of the pie chart but with just a different texture.

What I love about this metaphor is its instant clarity. That small piece calls out, “I’m small, but mighty!”

P.S. When you’re ready to learn more frameworks like this, you can \textbf{secure your spot for Thinking in Visual Metaphors. }Enrollment for Cohort 11 launches today!
\end{minipage}
\end{center}
\vspace*{\fill}